% !TeX root = ../main.tex
% =================================================================================
\chapter{Evaluation}\label{chap:evaluation}
% =================================================================================

This chapter evaluates our approach (contribution \refC{evaluation}).
First, \cref{sec:evaluation:gqm} introduces the evaluation plan using the \gls{gqm} approach.
\Cref{sec:evaluation:design} describes the study design, and \cref{sec:evaluation:results} presents the results.
\Cref{sec:evaluation:discussion} discusses the results with respect to our research questions, and \cref{sec:evaluation:threats} discusses threats to validity.

% ---------------------------------------------------------------------------------
\section{Goal-Question-Metric Plan}\label{sec:evaluation:gqm}
% ---------------------------------------------------------------------------------

The evaluation is structured according to the \gls{gqm} approach~\cite{Basili1994}.
Each goal is broken down into questions that tell whether the goal is reached, and each question into metrics that answer it.
Derive the goals from your research questions and requirements, for example:

\begin{description}[leftmargin=2.5em, labelindent=0pt, font=\normalfont]
  \item[$\boldsymbol{G}_1$] Improve the usability of \ldots for \ldots
  \begin{description}[leftmargin=2.5em, labelindent=0pt, font=\normalfont]
    \item[$\boldsymbol{Q}_{1,1}$] Does our approach achieve a higher usability than the baseline?
    \begin{description}[leftmargin=3em, labelindent=0pt, font=\normalfont]
      \item[$\boldsymbol{M}_{1,1,1}$] The System Usability Scale score (0--100) of both tools, collected from $n$ participants.
            Hypothesis $H_{A,1}$: the score of our approach is higher than the score of the baseline.
    \end{description}
    \item[$\boldsymbol{Q}_{1,2}$] Do users complete tasks faster with our approach?
    \begin{description}[leftmargin=3em, labelindent=0pt, font=\normalfont]
      \item[$\boldsymbol{M}_{1,2,1}$] Task completion time in minutes.
    \end{description}
  \end{description}
  \item[$\boldsymbol{G}_2$] \ldots
\end{description}

% ---------------------------------------------------------------------------------
\section{Study Design}\label{sec:evaluation:design}
% ---------------------------------------------------------------------------------

Describe the setup of your evaluation so that others can reproduce it:
the case studies, data, or tasks you use; the participants and how you recruited them; the procedure; the hardware and software; and how you analyze the data (e.g., which statistical tests).
Publish your data and scripts as a replication package and reference it here.

% ---------------------------------------------------------------------------------
\section{Results}\label{sec:evaluation:results}
% ---------------------------------------------------------------------------------

Report the results for each metric without interpreting them yet.
Reference every table and figure in the text and describe what it shows, e.g., \enquote{\Cref{tab:evaluation:results} shows the runtime of both approaches.}
Move extensive details to the appendix.

\begin{table}[htbp]
  \centering
  \caption{Runtime of the approaches (example data)}
  \label{tab:evaluation:results}
  \begin{tabular}{@{}l S[table-format=3.1] S[table-format=2.1] S[table-format=1.2]@{}}
    \toprule
    Approach & {Mean runtime (\unit{\milli\second})} & {SD (\unit{\milli\second})} & {Speedup} \\
    \midrule
    Baseline      & 120.4 & 12.3 & 1.00 \\
    Our approach  &  48.7 &  5.1 & 2.47 \\
    \bottomrule
  \end{tabular}
\end{table}

% ---------------------------------------------------------------------------------
\section{Discussion}\label{sec:evaluation:discussion}
% ---------------------------------------------------------------------------------

In this section, we discuss our research questions, taking the results of all parts of the evaluation into account.
Answer each research question explicitly, based on the metrics of \cref{sec:evaluation:gqm}, and summarize the key insights as findings.

\discussRQ{detect}

Discuss the results that answer the first research question and refer to the metrics, e.g., $\boldsymbol{M}_{1,1,1}$.
Interpret them: What do they mean, and why did you obtain them?

\finding[finding:faster]{Our approach detects modifications 2.5 times faster than the baseline.}

Continue the discussion and summarize further insights as findings, which you can reference later, e.g., \cref{finding:faster}.

\discussRQ{resolve}

Proceed in the same way for every research question.

\finding{Users resolve conflicts with our approach in less time than with the baseline.}

% ---------------------------------------------------------------------------------
\section{Threats to Validity}\label{sec:evaluation:threats}
% ---------------------------------------------------------------------------------

Discuss what limits the validity of your results and how you mitigated it:
\begin{description}
  \item[Construct validity] Do your metrics capture what you intended to evaluate?
  \item[Internal validity] Are the results caused by your approach, or by other factors?
  \item[External validity] To which other settings, data, or participants do the results generalize?
  \item[Reliability] Would other researchers obtain the same results with your replication package?
\end{description}
\textcite{Runeson2009} provide a checklist for such threats in software engineering case studies.
